Skill extraction is a prerequisite for systems that compare a person's evidence-backed capabilities with job requirements and learning opportunities. Yet a useful extraction system must do more than emit plausible labels: it must preserve the source span, distinguish hard and soft skills, remain robust when auxiliary components disagree, and expose enough provenance for later matching and recommendation. This report presents SkillLLM\_v2, a modular research system for span-level skill extraction, canonicalization, profile construction, graph-structured explanation, job matching, and course recommendation.

The central extraction component is a Mistral Small 3.1 24B Instruct model adapted with a low-rank adapter on the local SkillSpan export. The audited 80-epoch, two-GPU run obtains exact-span F1 of 0.670069, token F1 of 0.787886, hard-skill F1 of 0.726450, and soft-skill F1 of 0.604046 on its persisted full-test artifact. In the same evaluation record, the BERT and zero-shot GLiNER baselines achieve exact-span F1 of 0.059150 and 0.075833 respectively. These results make the adapted Mistral model the extraction anchor rather than evidence that adding more modules automatically improves performance.

SkillLLM\_v2 therefore uses provenance-aware, non-regressive composition. Candidate merging protects anchor spans, taxonomy linking is attach-only, and a gate selects wrapper predictions only when they stay within configured exact-span and soft-skill tolerances. The full wrapper is accepted as non-regressive and records exact-span F1 of 0.671102, but the difference is small and does not establish a general improvement claim. The implementation also constructs evidence-preserving profile summaries, deterministic cached context packets, and an in-memory skill graph. These downstream paths are implemented and smoke-tested; however, labelled job-ranking and course-ranking evaluations are unavailable, so no effectiveness claim is made for them.

The report contributes a detailed code-to-artifact audit, a transparent experimental account, reproducible system diagrams, and a research finding of practical importance: a modular career-intelligence pipeline should treat high-quality extraction as an invariant to protect, while using normalization, context, and graph components to enrich provenance and explanation rather than silently rewriting reliable spans.
